\section{Index notation}
The same ideas extend naturally from three dimensions to $\mathbb{R}^n$. Index notation provides a compact language for vectors, matrices, and tensors. An $n$-dimensional vector is an ordered list of components,
\begin{equation}
    \mathbf{v}=(v_1,v_2,\ldots,v_n)
    =\sum_{i=1}^{n}v_i\mathbf{e}_i,
\end{equation}
where $\{\mathbf{e}_1,\ldots,\mathbf{e}_n\}$ is an orthonormal Cartesian basis. Two vectors are equal if and only if all corresponding components are equal.

In $\mathbb{R}^n$, the coordinates of a point are $(x_1,x_2,\ldots,x_n)$, and a Cartesian basis is written as $\mathbf{e}_1,\mathbf{e}_2,\ldots,\mathbf{e}_n$. In index notation, the same vector is written
\begin{equation}
\mathbf{v}=\sum_{i=1}^{n}v_i\mathbf{e}_i.
\end{equation}
The entry in row $i$ and column $j$ of a matrix $M$ is denoted by $M_{ij}$.

\subsection{Einstein summation convention}
The Einstein summation convention suppresses an explicit summation sign: an index repeated exactly twice in a term is summed over its full range. Thus
\begin{equation}
\mathbf{u}\cdot\mathbf{v}=\sum_{i=1}^{n}u_iv_i=u_iv_i,
\end{equation}
and matrix multiplication becomes
\begin{equation}
[AB]_{ij}=\sum_{k=1}^{n}A_{ik}B_{kj}=A_{ik}B_{kj}.
\end{equation}
An index that is summed over is a \emph{dummy index}; an index that occurs once in each term is a \emph{free index}. Dummy indices may be renamed, while free indices must remain unchanged. For example,
\begin{equation}
A_{ik}B_{kl}C_{lj}=A_{im}B_{ml}C_{lj},
\end{equation}
but $A_{ik}B_{kk}C_{kj}$ is invalid because $k$ occurs three times. As a practical rule, each index should occur once or twice in any term, never more.

\subsection{Kronecker delta and Levi-Civita symbol}
The Kronecker delta is the component form of the identity matrix:
\begin{equation}
\delta_{ij}=\begin{cases}
1, & i=j,\\
0, & i\ne j.
\end{cases}
\end{equation}
Its useful contraction identities are
\begin{equation}
\delta_{ij}a_j=a_i,
\qquad
\delta_{ij}\delta_{jk}=\delta_{ik},
\qquad
\delta_{ii}=n.
\end{equation}
Proofs of these contractions are collected in Exercise~\ref{ex:index-delta}.

In three dimensions, the completely antisymmetric Levi-Civita symbol is
\begin{equation}
\varepsilon_{ijk}=\begin{cases}
+1, & (i,j,k)\text{ is an even permutation of }(1,2,3),\\
-1, & (i,j,k)\text{ is an odd permutation of }(1,2,3),\\
0, & \text{otherwise}.
\end{cases}
\end{equation}
For example, $\varepsilon_{123}=\varepsilon_{231}=\varepsilon_{312}=1$, while $\varepsilon_{132}=-1$. Interchanging any two indices changes its sign:
\begin{equation}
\varepsilon_{ijk}=-\varepsilon_{jik}.
\end{equation}
The antisymmetry argument and its consequence $\mathbf{u}\times\mathbf{u}=\mathbf{0}$ are developed in Exercise~\ref{ex:index-epsilon}.

\subsection{Vectors and matrices in index form}
The linear system $A\mathbf{x}=\mathbf{b}$ is written as
\begin{equation}
A_{ij}x_j=b_i.
\end{equation}
Similarly, the eigenvalue equation is
\begin{equation}
A_{ij}v_j^{(\alpha)}=\lambda^{(\alpha)}v_i^{(\alpha)},
\qquad \alpha=1,\ldots,n,
\end{equation}
where $\alpha$ labels an eigenpair and is not summed. An orthogonal matrix satisfies
\begin{equation}
R_{ki}R_{kj}=\delta_{ij},
\end{equation}
and its trace is $\operatorname{tr}A=A_{ii}$. The cyclic property of the trace follows directly from commutativity of scalar components:
\begin{equation}
\operatorname{tr}(AB)=A_{ij}B_{ji}=B_{ji}A_{ij}=\operatorname{tr}(BA).
\end{equation}

\begin{example}[Two-dimensional rotation]
For
\begin{equation}
R=\begin{pmatrix}\cos\theta&-\sin\theta\\ \sin\theta&\cos\theta\end{pmatrix},
\end{equation}
the $11$ component of $R^{\mathsf T}R$ is
\begin{equation}
R_{k1}R_{k1}=\cos^2\theta+\sin^2\theta=1,
\end{equation}
whereas the $12$ component is
\begin{equation}
R_{k1}R_{k2}=\cos\theta(-\sin\theta)+\sin\theta\cos\theta=0.
\end{equation}
Thus $R_{ki}R_{kj}=\delta_{ij}$.
\end{example}
The general preservation of inner products by orthogonal transformations is proved in Exercise~\ref{ex:index-orthogonal}.

\subsection{Dot and cross products}
For vectors with components $u_i$ and $v_i$,
\begin{equation}
\mathbf{u}\cdot\mathbf{v}=u_iv_i=\delta_{ij}u_iv_j,
\end{equation}
while the cross product is expressed compactly by
\begin{equation}
(\mathbf{u}\times\mathbf{v})_i=\varepsilon_{ijk}u_jv_k.
\end{equation}
Antisymmetry immediately gives $\mathbf{u}\times\mathbf{u}=\mathbf{0}$, since $u_ju_k$ is symmetric in $j,k$ whereas $\varepsilon_{ijk}$ is antisymmetric.
See Exercise~\ref{ex:index-epsilon} for the index-level proof.

\subsection{Triple products and epsilon-delta identities}
The scalar triple product has the index form
\begin{equation}
\mathbf{u}\cdot(\mathbf{v}\times\mathbf{w})=\varepsilon_{ijk}u_iv_jw_k.
\end{equation}
The vector triple product follows from the contraction identity
\begin{equation}
\varepsilon_{ijk}\varepsilon_{klm}=\delta_{il}\delta_{jm}-\delta_{im}\delta_{jl}:
\end{equation}
\begin{equation}
\mathbf{u}\times(\mathbf{v}\times\mathbf{w})
=(\mathbf{u}\cdot\mathbf{w})\mathbf{v}-(\mathbf{u}\cdot\mathbf{v})\mathbf{w}.
\end{equation}
Further useful contractions are
\begin{equation}
\varepsilon_{ijk}\varepsilon_{ljk}=2\delta_{il},
\qquad
\varepsilon_{ijk}\varepsilon_{ijk}=6,
\qquad
\varepsilon_{ijk}\delta_{jk}=0.
\end{equation}
The contraction identity and the vector triple-product derivation are left to Exercise~\ref{ex:index-triple}.

\subsection{Derivatives in index notation}
Write the partial derivative with respect to $x_i$ as
\begin{equation}
\frac{\partial}{\partial x_i}\longrightarrow\partial_i.
\end{equation}
For example, the vacuum Maxwell equations become
\begin{equation}
\begin{aligned}
\nabla\cdot\mathbf{E}=0
&\longrightarrow \partial_iE_i=0,
&\qquad
\nabla\cdot\mathbf{B}=0
&\longrightarrow \partial_iB_i=0,\\
\nabla\times\mathbf{E}+\partial_t\mathbf{B}=\mathbf{0}
&\longrightarrow \varepsilon_{ijk}\partial_jE_k+\partial_tB_i=0,
&\qquad
\nabla\times\mathbf{B}-\mu_0\varepsilon_0\partial_t\mathbf{E}=\mathbf{0}
&\longrightarrow \varepsilon_{ijk}\partial_jB_k-\mu_0\varepsilon_0\partial_tE_i=0.
\end{aligned}
\end{equation}

As an example, consider the curl of a gradient. In index notation,
\begin{equation}
\begin{aligned}
(\nabla\times\nabla\varphi)_i
&=\varepsilon_{ijk}\partial_j\partial_k\varphi\\
&=\varepsilon_{ijk}\partial_k\partial_j\varphi
&&\text{(mixed partials commute)}\\
&=-\varepsilon_{ikj}\partial_k\partial_j\varphi
&&\text{($\varepsilon_{ijk}$ is antisymmetric)}\\
&=-(\nabla\times\nabla\varphi)_i.
\end{aligned}
\end{equation}
Therefore,
\begin{equation}
\nabla\times\nabla\varphi=\mathbf{0}.
\end{equation}
See Exercise~\ref{ex:index-curl-gradient} for this proof in exercise form.

The curl-curl identity follows by applying the epsilon-delta contraction identity:
\begin{equation}
\begin{aligned}
\bigl[\nabla\times(\nabla\times\mathbf{E})\bigr]_i
&=\varepsilon_{ijk}\partial_j
    \left(\varepsilon_{klm}\partial_lE_m\right)\\
&=(\delta_{il}\delta_{jm}-\delta_{im}\delta_{jl})
    \partial_j\partial_lE_m\\
&=\partial_i\partial_jE_j-\partial_j\partial_jE_i.
\end{aligned}
\end{equation}
Returning to vector notation gives
\begin{equation}
\nabla\times(\nabla\times\mathbf{E})
=\nabla(\nabla\cdot\mathbf{E})-\nabla^2\mathbf{E}.
\end{equation}
See Exercise~\ref{ex:index-curl-curl} for a proof using index notation.

\subsection{Tensors and coordinate transformations}
A rank-two tensor maps one vector to another linearly. For example, stress and conductivity take the forms
\begin{equation}
f_i=\sigma_{ij}\,\mathrm{d}A_j,
\qquad
J_i=\sigma_{ij}E_j.
\end{equation}
Under an orthogonal coordinate transformation $x'_i=R_{ij}x_j$, vector components transform as $v'_i=R_{ij}v_j$, while a rank-two tensor transforms as
\begin{equation}
T'_{ij}=R_{ik}R_{jl}T_{kl}.
\end{equation}
More generally, a rank-$n$ tensor transforms according to
\begin{equation}
T'_{i_1\cdots i_n}
=R_{i_1j_1}\cdots R_{i_nj_n}T_{j_1\cdots j_n}.
\end{equation}

An isotropic tensor has the same components in every orthogonal coordinate system. The Kronecker delta is isotropic, since
\begin{equation}
R_{ik}R_{jl}\delta_{kl}=\delta_{ij}.
\end{equation}

